\section{Limitations and Future Work}
\label{sec:limitations}
%%%%%%%%%%%%%%%%%%%%%%%%%%%%%%%%%%%%%%%%%%%%%%%%%%%%%%%%%%%%%%%%%%%%%%%%%%%%%%%

Every theoretical framework involves tradeoffs between generality, precision, and realism.
This section provides an honest assessment of where the $(C, \Psi)$ framework
falls short, what it cannot do, and where future work is needed.

\subsection{Theoretical Limitations}

\subsubsection{The Contraction Assumption}

\paragraph{The Problem.}
The existence and uniqueness of $C^*$ (Theorem~\ref{thm:existence}) depends on 
the behavioral response operator $\Phi$ being a contraction:
\begin{equation}
  \|\Phi(C_1, \Psi) - \Phi(C_2, \Psi)\|_F \leq \gamma \|C_1 - C_2\|_F, \quad \gamma < 1
\end{equation}

This assumption \emph{fails} precisely when the framework is most needed:
\begin{itemize}
  \item Financial panics (beliefs about others' beliefs spiral)
  \item Bank runs (withdrawal expectations self-fulfill explosively)
  \item Revolutions (coordination on regime change)
  \item Hyperinflation (price expectations accelerate)
\end{itemize}

\paragraph{What This Means.}
The framework describes stability and gradual change well.
It has limited purchase on \emph{explosive} dynamics---
the moments when systems jump between equilibria.

\paragraph{Possible Resolution.}
Future work could:
\begin{itemize}
  \item Develop regime-switching models where $\gamma < 1$ in normal times, $\gamma > 1$ during crises
  \item Use catastrophe theory to model discontinuous transitions
  \item Integrate agent-based modeling \citep{tesfatsion2006} for micro-foundations of explosive dynamics
\end{itemize}

\subsubsection{Linearity of Context Dynamics}

\paragraph{The Problem.}
The context update equation
\begin{equation}
  \Psi_{t+1} = \Psi_t + \eta \cdot \Delta a_t
\end{equation}
is a first-order linear approximation. Real institutional change exhibits:
\begin{itemize}
  \item \textbf{Thresholds:} Change only after critical mass (S-curve dynamics)
  \item \textbf{Hysteresis:} Path back differs from path forward
  \item \textbf{Irreversibility:} Some changes cannot be undone (extinctions, lost knowledge)
  \item \textbf{Discontinuities:} Sudden jumps (revolutions, technological breakthroughs)
\end{itemize}

\paragraph{What This Means.}
The linear model captures tendencies but misses the nonlinear phenomena
that often matter most for policy.

\paragraph{Possible Resolution.}
Section~\ref{sec:context} sketched nonlinear extensions (S-curve, punctuated equilibrium).
Future work should:
\begin{itemize}
  \item Specify functional forms for different context types
  \item Empirically estimate adjustment speeds and threshold parameters
  \item Develop simulation tools for nonlinear context dynamics
\end{itemize}

\subsubsection{The Aggregation Problem}

\paragraph{The Problem.}
The welfare function aggregates across agents:
\begin{equation}
  Q = \sum_i \sum_d \sum_t \omega_i \delta_d^t [G_{i,d,t} - \lambda_d P_{i,d,t}]
\end{equation}

This utilitarian aggregation obscures:
\begin{itemize}
  \item \textbf{Distribution:} Two equilibria with identical $Q$ may differ dramatically 
        in how welfare is distributed
  \item \textbf{Power:} Who determines the weights $\omega_i$?
  \item \textbf{Capabilities:} Are people able to convert resources into welfare equally?
\end{itemize}

\paragraph{What This Means.}
The framework can compare equilibria by aggregate welfare
but cannot resolve debates about justice, equality, or rights.

\paragraph{Possible Resolution.}
Future work could:
\begin{itemize}
  \item Incorporate Rawlsian maximin (maximize welfare of worst-off)
  \item Add distributional constraints ($Q$ subject to Gini $< \bar{G}$)
  \item Develop capability-adjusted welfare measures following Sen
  \item Separate efficiency from equity analysis
\end{itemize}

\subsubsection{No Theory of Equilibrium Selection}

\paragraph{The Problem.}
The framework identifies that multiple equilibria exist ($C^*_A$, $C^*_B$, ...)
but provides no theory of \emph{which} equilibrium will be selected.

Selection depends on:
\begin{itemize}
  \item History and initial conditions
  \item Focal points and coordination devices
  \item Power and political conflict
  \item Chance events and path dependence
\end{itemize}

None of these are modeled.

\paragraph{What This Means.}
The framework can say ``these are the possible coherent configurations''
but not ``this is the configuration that will emerge.''

\paragraph{Possible Resolution.}
Future work could:
\begin{itemize}
  \item Integrate evolutionary game theory for selection dynamics
  \item Add political economy models for institutional choice
  \item Develop historical case studies of equilibrium selection
  \item Use agent-based models to simulate selection processes
\end{itemize}

\subsection{Empirical Limitations}

\subsubsection{Identification of Complementarities}

\paragraph{The Problem.}
The complementarity matrix $C$ consists of second derivatives of \emph{unobserved}
utility functions:
\begin{equation}
  C_{ij} = \frac{\partial^2 U_i}{\partial x_i \partial x_j}
\end{equation}

Utility is not directly observable. Identification requires:
\begin{itemize}
  \item Strong functional form assumptions
  \item Exogenous variation in $x_j$ to estimate response of $x_i$
  \item Panel data to separate preference heterogeneity from complementarity
\end{itemize}

\paragraph{What This Means.}
Empirically estimating $C$ is extremely difficult.
Most empirical work will rely on proxies and reduced-form evidence.

\paragraph{Possible Resolution.}
\begin{itemize}
  \item Use behavioral experiments to estimate $C$ in controlled settings
  \item Exploit natural experiments (policy changes, shocks) for identification
  \item Develop structural estimation methods with credible instruments
  \item Accept that some applications will be qualitative rather than quantitative
\end{itemize}

\subsubsection{Parameter Proliferation}

\paragraph{The Problem.}
The full FEPSDE model has:
\begin{itemize}
  \item 144 utility components
  \item 6 dimension-specific discount rates $\delta_d$
  \item 6 dimension-specific loss aversion parameters $\lambda_d$
  \item 3 category weights $\omega_{INU}, \omega_{KNU}, \omega_{IDN}$
  \item Plus the full $144 \times 144$ complementarity matrix
\end{itemize}

This is an enormous parameter space.

\paragraph{What This Means.}
\begin{itemize}
  \item The full model cannot be estimated without heroic assumptions
  \item Risk of overfitting: with enough parameters, any data can be explained
  \item Different parameter choices can generate opposite predictions
\end{itemize}

\paragraph{Possible Resolution.}
\begin{itemize}
  \item Use simplified versions for specific applications (e.g., 2-dimension models)
  \item Impose theoretically motivated restrictions (e.g., symmetric $C$)
  \item Calibrate parameters from existing literature rather than estimating all
  \item Focus on qualitative predictions that are robust to parameter values
  \item Use Bayesian methods with informative priors
\end{itemize}

\subsubsection{No Empirical Validation Yet}

\paragraph{The Problem.}
This paper is theoretical. No empirical tests have been conducted.
The predictions in Section~\ref{sec:novelpredictions} remain untested.

\paragraph{What This Means.}
The framework's empirical relevance is asserted, not demonstrated.
It could be that real-world data reject the predictions.

\paragraph{Possible Resolution.}
Future work must:
\begin{itemize}
  \item Test cross-domain spillover predictions using financial/political data
  \item Estimate asymmetric recovery dynamics from crisis data
  \item Validate innovation-coherence tradeoff with technology adoption data
  \item Conduct experimental tests of identity-economics predictions
\end{itemize}

\subsubsection{WEIRD Bias}

\paragraph{The Problem.}
The empirical foundations (Section~\ref{sec:empiricalfoundations}) draw primarily on:
\begin{itemize}
  \item Laboratory experiments with university students
  \item Data from Western, Educated, Industrialized, Rich, Democratic populations
  \item English-language academic literature
\end{itemize}

\citet{henrich2010} showed that WEIRD populations are often outliers,
not representative of human psychology globally.

\paragraph{What This Means.}
\begin{itemize}
  \item Parameters like $\lambda$ (loss aversion) may vary cross-culturally
  \item The relative importance of FEPSDE dimensions may differ
  \item Collectivist societies may weight $KNU$ higher than individualist ones
  \item Some mechanisms may be culture-specific
\end{itemize}

\paragraph{Possible Resolution.}
\begin{itemize}
  \item Replicate key experiments in non-WEIRD populations
  \item Develop culture-specific parameterizations
  \item Test whether framework structure holds even if parameters differ
  \item Collaborate with researchers from diverse contexts
\end{itemize}

\subsection{Conceptual Limitations}

\subsubsection{Normative Neutrality}

\paragraph{The Problem.}
The framework ranks equilibria by aggregate FEPSDE welfare,
but the weights are not determined by the framework itself.

\begin{itemize}
  \item How much weight on Financial vs. Ecological?
  \item How much to discount the future?
  \item Whose utility counts (citizens, residents, future generations)?
  \item How to weight INU vs. KNU vs. IDN?
\end{itemize}

These are political and ethical choices, not scientific ones.

\paragraph{What This Means.}
The framework cannot resolve normative disagreements.
It can clarify what is at stake but not determine the right answer.

\paragraph{Our Position.}
This is a feature, not a bug. A scientific framework should not
dictate values. Its role is to:
\begin{itemize}
  \item Make tradeoffs explicit
  \item Show consequences of different weight choices
  \item Provide vocabulary for democratic deliberation
  \item Let societies choose their own weights
\end{itemize}

\subsubsection{Micro-Foundations}

\paragraph{The Problem.}
The framework operates at the system level (complementarity structures, coherence).
It does not provide micro-foundations for:
\begin{itemize}
  \item How individuals form expectations
  \item How agents learn about $C^*$
  \item How coordination happens in practice
  \item What cognitive processes underlie coherence-seeking
\end{itemize}

\paragraph{What This Means.}
The framework describes equilibrium properties but not the process
by which equilibria are reached. It is a ``as if'' model, not a
literal description of individual cognition.

\paragraph{Possible Resolution.}
\begin{itemize}
  \item Integrate with learning models (Bayesian updating, reinforcement learning)
  \item Develop agent-based micro-foundations
  \item Connect to cognitive science and neuroscience of social coordination
\end{itemize}

\subsubsection{Static vs. Dynamic Analysis}

\paragraph{The Problem.}
Most of the formal analysis is comparative statics:
comparing equilibria, not analyzing transitions between them.

But policy-relevant questions are often dynamic:
\begin{itemize}
  \item How long does transition take?
  \item What is the path of $K$ during transformation?
  \item When do systems get stuck in transition?
\end{itemize}

\paragraph{Possible Resolution.}
\begin{itemize}
  \item Develop explicit dynamic models of transition
  \item Simulate adjustment paths with different policies
  \item Study historical transitions to calibrate adjustment speeds
\end{itemize}

\subsection{What the Framework Cannot Do}

To be explicit about scope:

\begin{enumerate}
  \item \textbf{Cannot predict which equilibrium will be selected.}
        It identifies possible equilibria but not which will emerge.
  
  \item \textbf{Cannot resolve normative disagreements.}
        It clarifies tradeoffs but cannot say what is right.
  
  \item \textbf{Cannot model explosive dynamics.}
        It describes stability, not the moments when stability breaks down.
  
  \item \textbf{Cannot explain individual behavior in detail.}
        It models system-level coherence, not individual psychology.
  
  \item \textbf{Cannot be fully empirically validated with current data.}
        Many predictions require data that doesn't yet exist.
  
  \item \textbf{Cannot capture everything.}
        144 components is already complex; reality is more complex still.
\end{enumerate}

\subsection{Epistemic Honesty: Integration vs. Prediction}

A critical distinction must be made explicit.

\paragraph{The Honest Claim.}
This framework is an \emph{integrative schema}, not a \emph{predictive theory}
that anticipated the contributions it encompasses.

Appendix~\ref{app:nobel} shows how 56 years of Nobel Prize research maps onto
the framework. But this mapping was constructed \emph{ex post}---after the
contributions were made. The framework did not predict Kahneman's loss aversion,
Ostrom's polycentric governance, or Romer's endogenous growth. These were
genuine scientific discoveries that \emph{extended} economic understanding.

\paragraph{What Would Be Dishonest.}
It would be intellectually dishonest to claim that:
\begin{itemize}
  \item The framework ``contains'' Prospect Theory (Kahneman 2002)
  \item The framework ``anticipated'' commons governance (Ostrom 2009)
  \item The framework ``predicted'' endogenous growth (Romer 2018)
  \item The framework ``implied'' institutional causality (Acemoglu 2024)
\end{itemize}

Each of these was a genuine discovery that required creativity, empirical work,
and theoretical innovation. The framework was \emph{built to accommodate them},
not the reverse.

\paragraph{Genuine Extensions in the Nobel Record.}
Several Nobel contributions represent \emph{true extensions} that expanded
economic understanding in ways no prior framework anticipated:

\begin{center}
\renewcommand{\arraystretch}{1.3}
\small
\begin{tabular}{lll}
\hline
\textbf{Laureate} & \textbf{Discovery} & \textbf{Framework Adaptation} \\
\hline
Kahneman (2002) & Loss aversion, reference points & Added $\lambda_d > 1$, $\Psi_{frame}$ \\
Simon (1978) & Bounded rationality & $C^{perceived} \neq C^{actual}$ \\
Ostrom (2009) & Polycentric governance & $\Psi_{inst}^{community}$ as third way \\
Thaler (2017) & Choice architecture & $\Psi_{architecture}$ as policy tool \\
Romer (2018) & Endogenous technology & $\Psi_{tech,t+1} = f(\Psi_{tech,t}, R\&D)$ \\
AJR (2024) & Causal identification & Empirical test of $\Psi_{inst} \to Q$ \\
\hline
\end{tabular}
\end{center}

These were not ``special cases'' of a pre-existing theory.
They were \emph{expansions of knowledge} that the framework subsequently incorporated.

\paragraph{What the Framework Actually Does.}
The honest characterization is that the framework:
\begin{enumerate}
  \item \textbf{Provides common language:} Different contributions can be 
        compared using $(C, \Psi, C^*, K, Q)$
  \item \textbf{Makes assumptions explicit:} Each theory's restrictions become visible
  \item \textbf{Identifies connections:} Links between disparate literatures emerge
  \item \textbf{Reveals gaps:} What has \emph{not} been studied becomes clear
  \item \textbf{Suggests extensions:} Where might the next discoveries lie?
\end{enumerate}

\paragraph{What the Framework Does NOT Do.}
\begin{enumerate}
  \item \textbf{Does not replace creativity:} Genuine discoveries require insight
        beyond any schema
  \item \textbf{Does not predict future Nobels:} We cannot say what will be 
        discovered next
  \item \textbf{Does not make prior work redundant:} Kahneman's experiments,
        Ostrom's fieldwork, Romer's mathematics remain essential
  \item \textbf{Does not guarantee correctness:} The framework may need revision
        when future discoveries don't fit
\end{enumerate}

\paragraph{The Retrospective Fitting Problem.}
Any sufficiently flexible framework can accommodate past findings.
The test is whether it generates \emph{novel, falsifiable predictions}
that would not have been made without it.

Section~\ref{sec:predictions} attempts this. But honesty requires acknowledging:
\begin{itemize}
  \item The predictions have not yet been tested
  \item Some predictions may be too vague to falsify
  \item Future research may reveal predictions that fail
\end{itemize}

\paragraph{Why This Matters.}
Science progresses through genuine discovery, not retrospective categorization.
The framework's value lies in facilitating communication, identifying gaps,
and suggesting directions---not in claiming to have ``known all along.''

The 56 Nobel laureates in Appendix~\ref{app:nobel} made genuine contributions.
The framework honors them by providing a structure for understanding how their
insights connect---not by claiming credit for their creativity.

\paragraph{A Constructive Stance.}
Rather than asking ``Does this framework contain all prior knowledge?''
(to which the honest answer is ``no---it was built to integrate it''),
the productive question is:

\begin{quote}
\emph{``What does this framework suggest we should study next?''}
\end{quote}

Some candidates:
\begin{itemize}
  \item Cross-level coherence conflicts (individual vs. collective optimization)
  \item Digital context ($\Psi_{digital}$) and its rapid evolution
  \item Ecological-economic complementarities under climate change
  \item Identity formation and its interaction with institutional change
\end{itemize}

These are not ``implied by'' the framework in any predictive sense.
They are \emph{suggested by} the framework's structure as potentially important.

Whether they yield genuine insights is an empirical question
that only future research can answer.

\paragraph{Retrospective vs. Prospective Application.}
A critical distinction clarifies the framework's epistemic status:

\textbf{Retrospective (1969--2025):} The framework was constructed \emph{after}
the Nobel contributions were made. Showing that Kahneman, Ostrom, or Romer
``fit'' the framework is \emph{integration}, not \emph{prediction}.
This is valuable for understanding connections but does not demonstrate
predictive power.

\textbf{Prospective (2026 onward):} The framework now exists as a complete
structure. Any new paper published tomorrow can be immediately analyzed:
\begin{itemize}
  \item Which $C_{ij}$ does it address?
  \item Which $\Psi$ components are relevant?
  \item What restrictions does it assume?
  \item What would the full framework add?
  \item How does it connect to existing literature?
\end{itemize}

This prospective capability is the framework's \emph{operational value}.
It provides a systematic protocol for integrating new research as it appears.

\textbf{The Difference:}
\begin{center}
\renewcommand{\arraystretch}{1.3}
\begin{tabular}{lll}
\hline
& \textbf{Retrospective} & \textbf{Prospective} \\
\hline
Claim & ``Framework integrates X'' & ``Framework classifies X'' \\
Epistemic status & Post-hoc organization & Real-time analysis \\
Value & Understanding connections & Operational tool \\
Honesty required & ``Built to fit'' & ``Applies to new cases'' \\
\hline
\end{tabular}
\end{center}

The 56 Nobel laureates in Appendix~\ref{app:nobel} demonstrate retrospective
integration. The framework's prospective value will be demonstrated by
its ability to usefully classify research that does not yet exist.

\paragraph{A Testable Claim.}
Here is a falsifiable prediction about the framework itself:

\begin{quote}
\emph{Any economics paper published in 2026 or later can be characterized
in terms of $(C, \Psi, C^*, K, Q)$, with explicit identification of
which restrictions it assumes and which components it addresses.}
\end{quote}

If a paper appears that genuinely cannot be characterized this way---
not because it extends the framework (which is valuable) but because
it contradicts the framework's basic structure---then the framework
requires revision.

This is the standard any unified framework must meet:
not that it predicted everything, but that it can integrate anything---
or honestly acknowledge when it cannot.

\subsection{What the Framework Can Do}

Despite these limitations:

\begin{enumerate}
  \item \textbf{Provides unified language.}
        Economists, political scientists, psychologists, and organizational scholars
        can use $(C, \Psi, K, Q)$ to communicate across disciplines.
  
  \item \textbf{Generates novel predictions.}
        Cross-domain spillovers, asymmetric recovery, identity-economics feedback---
        these predictions are unavailable from constituent theories.
  
  \item \textbf{Clarifies policy tradeoffs.}
        The Trump tariff analysis shows how the framework illuminates
        tradeoffs that standard models miss.
  
  \item \textbf{Explains puzzles.}
        Why do people support economically costly policies?
        Why do stable systems suddenly collapse?
        Why is reform so hard? The framework offers answers.
  
  \item \textbf{Structures empirical research.}
        Even without full estimation, the framework suggests
        what to measure and what relationships to test.
  
  \item \textbf{Enables better questions.}
        ``What is the coherence index?'' is more productive than
        ``Is this policy efficient?'' in many contexts.
\end{enumerate}

\subsection{Research Agenda}

\paragraph{Theoretical Extensions.}
\begin{itemize}
  \item Nonlinear context dynamics with tipping points
  \item Equilibrium selection theory
  \item Dynamic transition analysis
  \item Multi-level coherence interactions
\end{itemize}

\paragraph{Empirical Program.}
\begin{itemize}
  \item Experimental estimation of complementarity parameters
  \item Cross-country validation of FEPSDE welfare measures
  \item Time-series tests of coherence spillover predictions
  \item Case studies of major transitions (EU integration, China's reform)
\end{itemize}

\paragraph{Applied Work.}
\begin{itemize}
  \item Trade policy analysis (beyond tariffs: supply chain resilience)
  \item Climate policy (global coherence problem)
  \item Digital transformation (innovation-coherence tradeoff)
  \item Institutional reform (coherence trap escape)
\end{itemize}

\paragraph{Methodological Development.}
\begin{itemize}
  \item Simulation tools for coherence dynamics
  \item Estimation methods for high-dimensional complementarity structures
  \item Survey instruments for FEPSDE measurement
  \item Case study protocols for qualitative application
\end{itemize}

The feasibility of this research agenda depends critically on
the computational infrastructure described throughout this paper.
The 60-year evolution from manual statistical analysis to AI-assisted
economic modeling---documented in Appendix~\ref{app:computationalhistory}---is
what transforms these methodological goals from aspirations to achievable objectives.

\subsection{Conclusion: The Value of Imperfect Frameworks}

Every framework has limitations. The question is not whether limitations exist
but whether the framework is useful despite them.

The $(C, \Psi)$ framework is useful because it:
\begin{itemize}
  \item Asks questions that other frameworks don't ask
  \item Generates predictions that other frameworks don't generate
  \item Connects literatures that other frameworks don't connect
  \item Provides vocabulary that other frameworks don't provide
\end{itemize}

Its limitations are acknowledged. Its value lies in providing
a coherent structure for asking better questions about
stability, welfare, and institutional design---
even when definitive answers remain elusive.

%%%%%%%%%%%%%%%%%%%%%%%%%%%%%%%%%%%%%%%%%%%%%%%%%%%%%%%%%%%%%%%%%%%%%%%%%%%%%%%